\chapter{风险的价格：为什么有人借钱 3\%，有人 12\%}
\chaptermeta{13}

\begin{ledetext}
管道已经铺好，钱能流到任何地方。但流过去要收多少过路费，是另一门手艺。这门手艺一半是算术，一半是人心。
\end{ledetext}

\section{郑姐那天下午问了四个地方}

郑姐在常州开一家做汽车密封条的小厂，二十三个工人，2025 年出货一千一百多万。2026 年 8 月的一个下午，她要凑 80 万买一台新的注塑机。旧的那台漏油漏了大半年，修一次两万三，一年修了四回。

她问了四个地方，拿回四个数字。

\begin{tablewrap}
\begin{xltabular}{\linewidth}{>{\raggedright\arraybackslash}p{0.20\linewidth}XXr}
\toprule
\bthead{路子} & \bthead{拿什么做保} & \bthead{条件} & \bthead{年化} \\
\midrule
\textbf{房贷（她家那套 2019 年买的）} & 房子本身 & 30 年，等额本息 & 3.5\% \\
\textbf{抵押经营贷} & 住宅二次抵押 + 担保公司兜底 & 3 年，先息后本 & 3.9\% \\
\textbf{信用卡分期} & 什么都不用押 & 12 期，月费率 0.6\% & 13.03\% \\
\textbf{某持牌消费金融的小额经营贷} & 只看征信和流水 & 12 期，随借随还 & 23.9\% \\
\textbf{（对照）美国的私募信贷} & 企业资产 + 一堆财务契约 & 5 到 7 年，浮动 & 8\%–12\% \\
\bottomrule
\end{xltabular}
\end{tablewrap}

房贷那一行的 3.5\%，是 2026 年 8 月 20 日公布的 5 年期以上 LPR，已经连续多月没动。最后一行是给你看看，同一个世界的另一头，钱是什么价。

同一个人，同一天，同一笔 80 万。最贵的比最便宜的贵了将近七倍。

这不是谁比谁善良。

\begin{egbox}{先把那个 13.03\% 算出来}
客户经理跟郑姐说的是：月费率只有 0.6\%，一年下来 7.2\%，比经营贷贵一点点，好处是当天到账。

7.2\% 是假的。

她先试了 12 万，分 12 期。每月还本金 1 万，手续费按\textbf{最初的 12 万}算，120000 × 0.6\% = 720 元，所以每月还 10720。全年手续费 8640 元，8640 ÷ 120000 = 7.2\%，账面上没错。

错在分母。第一个月她手上确实压着 12 万，第二个月只剩 11 万，到第十二个月只剩 1 万。整段时间里她平均占用的钱是 6.5 万左右，不是 12 万。

拿 8640 除以 65000，13.3\%。严格解出内部收益率，是 \hlmark{13.03\%}。

银行说的每个字都对，它只是没提分母在缩水。中国从 2021 年起要求消费信贷明示年化利率，所以现在很多 App 上会有一行小字写着"年化利率（单利）"。看那一行。

\end{egbox}

\section{一笔利率是四块钱叠起来的}

任何一个利率数字，撕开都是同样几层。

\begin{flowlist}
  \flowitem{01}{无风险利率}{把钱借给几乎不可能赖账的对手，同样的期限，能拿多少。中国这边的锚是国债和政策利率，2026 年 8 月 7 天期逆回购利率是 1.4\%。这一块跟你是谁没关系，只跟时间有关。}
  \flowitem{02}{\term{期限溢价}{term premium}\index{042@期限溢价}}{借 3 年比借 7 天贵。放贷的人把钱锁死三年，这三年里通胀怎么走、政策怎么变、他自己会不会急用，全都没得反悔。这份不自由要收钱。}
  \flowitem{03}{\term{信用利差}{credit spread}\index{072@信用利差}}{借款人可能还不上。这是四块里最大、也最难算的一块，下一节整节都在算它。}
  \flowitem{04}{\term{流动性溢价}{liquidity premium}\index{035@流动性溢价}}{如果放贷方半路急需用钱，这笔债权能不能出手、要打几折。上市公司的债几分钟就能卖掉，郑姐这笔贷款卖给谁？}
  \flowitem{05}{税和条款}{抵押、担保、提前还款罚金、财务契约、要不要开一个受监管的资金账户。条款每收紧一格，利率就能松一格。这是可以谈的部分，也是小企业主最少去谈的部分。}
\end{flowlist}

教科书就分到这里为止。银行内部的定价表长得不太一样，它更实用，也更冷酷。

下面这个拆法是示意，不是哪家银行的真实报价单。假设这笔小微经营贷最后定在 8.1\%：

\begin{barfig}
  \barrow{资金成本}{25}{2.0\%}
  \barrow{预期损失}{30}{2.4\%}
  \barrow{运营成本}{27}{2.2\%}
  \barrow{资本回报}{15}{1.2\%}
  \barrow{税与杂项}{3}{0.3\%}
\end{barfig}

最长的那根不是银行的利润，是预期损失。

\section{2.4 个百分点，是九十六个人替四个人付的}

信贷定价的核心只有一个公式，写出来三个字母：

\begin{pullquote}{}
预期损失 = 违约概率 × 违约损失率 × 敞口
\end{pullquote}

这个式子朴素到有点无聊，但它决定了这个星球上每一笔贷款的价格。

\begin{egbox}{给一百家小微企业放贷}
一家城商行放 100 笔，每笔 80 万，总敞口 8000 万。

一年下来，4 家还不上。这是\term{违约概率}{probability of default}\index{063@违约概率}，4\%。

违约不等于全赔光。抵押品处置、诉讼、执行，平均两年半之后，每家收回 40\%。损失掉的是 60\%，这是\term{违约损失率}{loss given default}\index{064@违约损失率}。

算总账：4 笔 × 80 万 × 60\% = \textbf{192 万}。摊到 8000 万的盘子上，2.4\%。

也就是说，这家银行光要把本金亏掉的部分挣回来，利率就得比无风险利率高出 2.4 个点，一分钱不赚。

再加上资金成本 2.0\%、运营成本 2.2\%、股东要的资本回报 1.2\%，8\% 出头。\hlmark{一分利润没多拿。}

\end{egbox}

运营成本那 2.2\% 是最容易被忽略的一块。一笔 80 万的小微贷，客户经理要上门、要看厂房、要核订单、要查上下游，放款之后每季度还得跟踪，出问题要催收。这一整套动作，和一笔 8000 万的公司贷款做的事差不多。

人力成本几乎一样，分母差一百倍。

所以"小微企业融资难融资贵"这句话，前半段的第一层解释不是道德，是四家企业还不上，剩下九十六家在替他们付钱，外加一百套尽调的工时。

\section{但数学解释不了全部}

如果只有数学，郑姐的价格就该是 8\%，不该有 3.9\% 那一行。

3.9\% 是怎么来的？她把住宅做了二次抵押，还找了本地一家政府背景的融资担保公司。这两件事同时干了两件活：抵押品把违约损失率从 60\% 压到 25\% 左右，担保把银行实际承担的违约概率又切掉一大半。预期损失从 2.4\% 掉到不足 0.6\%，利率自然就下来了。

征信数据干的是另一半活。银行能合法看到她的增值税发票流水、社保缴纳、用电量、上下游的应收账款，那个 4\% 的违约概率就不再是一个行业平均数，而是她这一家的数。\hlwarm{风险定价贵，很多时候不是因为风险真的大，是因为看不清。}看不清的部分要按最坏的算。

我的判断是这样的：中国小微融资贵这件事，数学占大头，摩擦占的那一小半是可以往下压的，而且这些年确实压下去了一截。压它的从来不是文件里的道德呼吁，是税务数据、电力数据、供应链单据能不能被银行合法看见、看见之后能不能被采信。数据确权往前走一步，利率就往下挪一小格。这条路很慢，走了十几年才挪了这么点，但它是真的路。

\begin{mythbox}{常见误解}
\mvfrow{误解}{"利率收这么高，这家机构就是黑心。"}
\mvfrow{事实}{先把 8\% 和 24\% 分开看。8\% 里面绝大部分是预期损失和运营成本，是算出来的。24\% 那一档，预期损失可能真的有十几个点，因为来借的人本身就是被前面几道门槛筛剩下的。\\ 真正该骂的不是数字大，是\textbf{把 13.03\% 说成 7.2\%}，是把砍头息、服务费、保险费拆成七八个名目让你算不清总价。价格高不是罪，价格不透明才是。}
\end{mythbox}

\section{风险、不确定性，和一个消不掉的东西}

1921 年，芝加哥大学的弗兰克·奈特写了一本书，把两样长得很像的东西分开了。

骰子是\textbf{风险}：六个面，每一面的概率你都知道，你算不准下一把，但你算得准一千把。明天开门的那家新店能不能活过三年是\textbf{不确定性}：你不知道有几个面，甚至不知道它是不是个骰子。

金融业干的事，是把不确定性硬塞进风险的模子里。办法是找历史数据来充当那六个面。

数据有多长，你的骰子就有几个面。2007 年以前，美国全国房价从来没有同时下跌过，模型里根本不存在这一面。不是模型算错了，是那一面在骰子上没刻。

\begin{warnbox}{小心：VaR 那句话的后半句}
风控报表上最常见的数字叫\term{在险价值}{Value at Risk}\index{081@在险价值}。一句标准表述是："99\% 的置信度下，本组合单日亏损不超过 3200 万。"

这句话听着让人安心，因为它只说了 99\% 的日子。剩下那 1\% 的日子亏多少？\textbf{VaR 不回答。}它的定义里就不包含那部分。

更麻烦的是，那 1\% 的日子有个坏习惯：它们喜欢连着来。一年 250 个交易日里的两三天，往往不是散落在三月、七月和十一月，而是挤在同一个星期里。

\end{warnbox}

说完 1\% 的日子，说说另一件被误会得更深的事。

把钱摊到很多篮子里，能消掉的叫非系统性风险：这家厂子着火、那家公司的 CFO 跑路、某个专利官司输了。这些事互相不认识，一多就互相抵消。

消不掉的叫系统性风险：利率、汇率、战争、政策、全社会同时不敢花钱。它作用在所有篮子上。

决定你到底分散了多少的，只有一个变量：相关性。

\begin{egbox}{同样 20 只股票，三种命运}
假设你持有 20 只股票，每只 5\%，每只的年化波动都是 30\%。

如果它们两两之间完全不相关（相关系数 0），组合的波动会降到单只的 22.4\%，也就是 6.7\%。这就是分散化被称为"免费午餐"的原因，你没牺牲任何期望收益，风险少了四分之三。

把两两相关系数调到 0.6，同样这 20 只，组合波动只能降到单只的 \textbf{79\%}。

调到 0.9，降到 \textbf{95\%}。持有 20 只和持有 1 只，几乎没有区别。

2008 年 10 月和 2020 年 3 月，很多平时相关系数只有 0.2 的东西，在同一周里一起冲到了 0.9 以上。股票、公司债、新兴市场货币、黄金、REITs，全跌。

\end{egbox}

因为在恐慌里，所有人卖的不是某个资产，是"任何能卖掉的东西"。这时候资产之间的基本面关系失效了，剩下的只有一个共同因子：谁需要现金。

\begin{pullquote}{}
危机中所有相关性都趋近于 1。你精心搭的那个组合，恰好在最需要它的那一天不再是组合。
\end{pullquote}

\section{为什么好人会先离场}

1970 年，阿克洛夫写了一篇被三家期刊退稿的论文，讲二手车市场。买家分不清好车和坏车，只肯出一个平均价；好车车主觉得亏，不卖了；市场上剩下的车更差；买家再压价。最后好车全部退场，只剩柠檬。

把二手车换成借款人，这个循环一模一样。

\begin{chainflow}
\chainnode{分不清好坏}\chainarrow \chainnode{只肯按平均定价}\chainarrow \chainnode{好借款人嫌贵走人}\chainarrow \chainnode{池子里的人更差}\chainarrow \chainnode{只好再涨价}\chainarrow \chainnode{更好的人再走}
\end{chainflow}

这叫\term{逆向选择}{adverse selection}\index{040@逆向选择}。它发生在借钱之前。

借到钱之后还有第二个问题。郑姐拿这 80 万去买注塑机，银行放心。可如果她拿去做了别的，比如投给她表弟那个据说很赚钱的项目，银行是不知道的。钱一旦到手，用钱的人和出钱的人，利益就分岔了。这叫\term{道德风险}{moral hazard}\index{010@道德风险}。

金融业几百年来发明的那些看着很烦的东西，全是在治这两个病。抵押品治的是逆向选择（敢押房子的人对自己有信心），也治道德风险（乱来会失去房子）。分期放款治道德风险：钱按工程进度给，你没法一次性拿走。财务契约治的也是它：负债率超过某个数，银行有权提前收贷。征信记录治的是逆向选择，它把"你上次是什么样"变成了这次的价格。

评级机构本该是治逆向选择的最大一件工具。它的工作就是替所有人分辨好车和柠檬。

它的收费结构是：\textbf{被评的那家公司付钱。}

这个安排的问题不需要解释。2008 年，一大批贴着 AAA 标签的结构化产品变成了废纸，那是史上最贵的一次演示。危机之后确实改了不少：内部防火墙、方法论公开、监管备案、发行人不能挑评级师。但付钱的还是被评的人。

所以我对评级的态度是：当一份有用的、带着已知偏向的参考读，不当结论用。一家公司被降级往往说明市场早就知道了，评级只是最后一个承认的。

\section{2026 年：利差是最早的那支体温计}

\begin{statrow}{3}
  \statitem{2.3}{万亿美元}{全球私募信贷管理规模（2025 年）}
  \statitem{8–12}{\%}{私募信贷典型借款利率}
  \statitem{200–400}{bp}{比银行贷款高出的部分}
\end{statrow}

\begin{databox}{2026 数据}
私募信贷 2025 年全球管理规模约 2.3 万亿美元，2020 年以来翻了一倍，年化增速超过 17\%。借款利率通常在 8\% 到 12\%，比银行贷款高 200 到 400 个基点；二顺位和夹层这类更靠后的层级，12\% 到 16\%。

它的行业敞口里软件业约占 30\%，其中投向 AI 领域的约 2000 亿美元，占总资产 9\% 左右。

2025 年 10 月，美国两家中小银行 Zions Bancorporation 和 Western Alliance 先后出现风险事件。中金把私募信贷称作"2 万亿美元的灰犀牛"，认为风险可能在未来 12 到 24 个月逐步传导。这是一家机构的判断，不是已发生的事实。

\end{databox}

为什么这些利差数字值得盯着？因为信用市场的报价比股市诚实。

\begin{talkbox}
  \talkline{新闻}{"某中型企业的债券收益率较同期限国债的利差，从 320 个基点走阔至 610 个基点。"}
  \talkline{翻译}{三个月前，把钱借给这家公司比借给政府多赚 3.2 个点；现在要多赚 6.1 个点才有人肯借。市场对它违约概率的估计，翻了将近一倍。它的股价可能还在原地，因为股东在做梦，债权人不做梦。}
\end{talkbox}

股东买的是这家公司变成十倍的可能性，所以坏消息可以被"故事"抵消一阵子。债权人最多拿回本金加利息，上不封顶的东西跟他没关系，他只关心一件事：这家公司会不会死。

所以坏消息在信用市场里先有价格。2007 年 2 月，次级抵押贷款相关指数已经在往下掉，标普 500 还在往上爬，一直爬到 10 月才见顶。八个月的时间差。

\begin{mythbox}{常见误解}
\mvfrow{误解}{"高风险高收益。我年轻，能承受风险，所以我该去拿高收益。"}
\mvfrow{事实}{风险溢价是一个\textbf{期望值}，不是一份合同。它的完整意思是"平均而言、长期来看、在足够多的独立样本上"，多出来那几个点会兑现。你这一次不是样本，你这一次是一次抽样。\\ 更要命的是另一半：市场只补偿\emph{没法被分散掉的}那部分风险。你把全部身家压在一只股票上，承担了比市场高得多的波动，但其中大部分是可以被分散掉的非系统性风险，理论上一分钱溢价都不该给你。这叫承担了不被补偿的风险。\\ 还有第三种：某些资产的高利率不是给你的风险补偿，是给渠道的销售费用。这不构成投资建议，只是提醒你去看费率表。}
\end{mythbox}

\bookbreak

郑姐最后选了 3.9\% 那条路，把房子押了。她跟我说了一句话，我一直记着："押房子我心里踏实，因为我知道我肯定能还上。要是我自己都不敢押，那我凭什么让银行信我。"

这句话把这一章的道理讲完了。

\begin{takeaway}{本章一句话}
利率不是一个数字，是一份拆得开的账单：无风险利率 + 期限溢价 + 信用利差 + 流动性溢价，加上一堆条款和摩擦。

\begin{itemize}
  \item 预期损失 = 违约概率 × 违约损失率 × 敞口。小微贷款利率高，一大半是这道算术的结果。
  \item 另一小半是摩擦，可以靠抵押品、征信数据和担保压下去，但压不到零。
  \item 风险是可测的，不确定性不是。VaR 只告诉你 99\% 的日子，麻烦全在剩下那 1\%。
  \item 分散化的效果由相关性决定，而危机里所有相关性都往 1 跑。
  \item 信用利差是这台机器上最早报警的那支温度计，比股市早。
\end{itemize}

\end{takeaway}

\begin{bridgebox}
\textbf{下一章}定价错了，未必致命。真正把小麻烦变成灾难的，是另一样东西。它就藏在你家的房贷合同里，而且大多数人从没意识到自己正用着它。
\end{bridgebox}

